\chapter*{Resumen ejecutivo}
\addcontentsline{toc}{chapter}{Resumen ejecutivo}

\begin{summarybox}
Este documento presenta una propuesta técnica y comercial de Lunabotics para Alestis Puerto Real, orientada a mejorar la logística interna asociada a la línea del HTP (\textit{Horizontal Tail Plane}) del A320 y a las operaciones relacionadas con la sección 19.1. La solución propuesta no pretende que los robots fabriquen directamente estructuras aeronáuticas; su alcance realista es reducir desplazamientos improductivos, esperas por kits y herramientas, incidencias de trazabilidad y riesgos FOD (\textit{Foreign Object Debris}) mediante robots móviles autónomos.
\end{summarybox}

La propuesta recibe el nombre de \textbf{Sistema AMR-FOD-Kitting para línea HTP A320}. Se basa en una plataforma AMR omnidireccional, representable en CoppeliaSim mediante una base tipo KUKA YouBot, KUKA Omnirob o equivalente industrial, equipada con LiDAR, cámara RGB-D, lectura RFID/QR, encoders, IMU, escáner láser de seguridad, parada de emergencia, HMI y conexión WiFi industrial. La selección se justifica por el entorno interior de una nave aeronáutica, la necesidad de maniobrar cerca de estaciones y útiles de gran tamaño, y la conveniencia de realizar movimientos laterales sin reorientar el vehículo.

El sistema contribuye al objetivo de escalar desde una producción aproximada de 40 HTP al mes hacia 80 HTP al mes o más en un horizonte de dos años. La contribución esperada no es una duplicación automática de la producción, sino una mejora gradual de disponibilidad operativa: menos recorridos manuales, menor tiempo de espera, entregas coordinadas, registro de materiales y herramientas, inspección preventiva FOD y una base de datos de eventos que facilite la supervisión industrial.

La actividad se estructura según el enunciado: entrevista simulada con experto (20\%), oferta técnica (60\%) y oferta comercial (20\%). El presupuesto se presenta para tres escenarios de despliegue: piloto con 1 robot, expansión con 7 robots y despliegue extendido con 17 robots. Las cifras se apoyan en referencias comerciales públicas (MiR250, sensores, RFID, HMI, WiFi industrial y rangos 2026 de AMR) y en datasets auxiliares de supuestos; no constituyen una cotización vinculante. Para reforzar la defensa, el informe incorpora matriz de trazabilidad problema--requisito--solución, KPIs de aceptación, plan de simulación en CoppeliaSim, matriz de riesgos y presentación equivalente a PowerPoint en formato PPTX.
